\documentclass[11pt]{article}

\usepackage[T1]{fontenc}
\usepackage{lmodern}
\usepackage{amsmath,amssymb,amsthm}
\usepackage{booktabs}
\usepackage{geometry}
\usepackage[hidelinks]{hyperref}
\usepackage{microtype}
\geometry{margin=1in}
\hypersetup{
  pdftitle={The Connectivity-Code Number of C3 square C4 Is 16},
  pdfauthor={Xiaojian Zeng}
}

\newtheorem{theorem}{Theorem}
\newcommand{\F}{\mathbb F}
\newcommand{\cC}{\mathcal C}
\newcommand{\bits}[1]{\ensuremath{\mathtt{#1}}}

\title{The Connectivity-Code Number of $C_3\mathbin{\square}C_4$ Is $16$}
\author{Xiaojian Zeng\\
  \small PBC School of Finance, Tsinghua University, Beijing 100190, China\\
  \small \texttt{xiaojianzeng@foxmail.com}}
\date{}

\begin{document}
\maketitle

\begin{abstract}
For a finite graph $H$, a connectivity code is a family of edge subsets
whose pairwise symmetric differences are connected spanning subgraphs of
$H$.  We determine the exact maximum for the first fixed non-diagonal
product of cycles considered here:
\[
  m(C_3\mathbin{\square}C_4)=16.
\]
The degree-four restriction bound gives the upper bound for arbitrary,
possibly nonlinear families.  For the lower bound we give four explicit
24-bit generators.  A compact table supplies an ordered spanning tree in
each of the 15 nonzero generated words, and therefore certifies all 120
pairwise differences of the resulting 16-word linear code.
\end{abstract}

\section{Context and statement}

Let $H$ be a finite simple graph.  Following Alon~\cite{Alon2024}, a
\emph{connectivity code} for $H$ is a family $\cC\subseteq 2^{E(H)}$ such
that $(V(H),A\mathbin{\triangle}B)$ is connected for every two distinct
$A,B\in\cC$.  Write $m(H)$ for the maximum cardinality of such a family.
The standard cut-restriction argument gives
$m(H)\leq 2^{k'(H)}\leq 2^{\delta(H)}$~\cite{Alon2024}.

Bai, Gao, Ma, and Wu proved that
$m(C_m\mathbin{\square}C_n)\leq 16$ for $m,n\geq3$ and obtained equality
for $m=n=3$~\cite[Proposition~1.7]{BaiEtAl2024}.  Tian's subsequent work
constructs extremal connectivity codes in infinitely many regular graphs
of every degree at least four, and explicitly records determination of
products of cycles as a remaining general problem~\cite[Section~6(d)]{Tian2026}.
The result below concerns only the single graph $C_3\mathbin{\square}C_4$;
it makes no claim about all products of cycles or about absolute priority
beyond those sources.

\begin{theorem}\label{thm:main}
Let $H=C_3\mathbin{\square}C_4$ on
$\mathbb Z_3\times\mathbb Z_4$.  If $m(H)$ is maximized over all families
of edge subsets, not only linear families, then
\[
  m(H)=16.
\]
\end{theorem}

\section{Upper bound and explicit construction}

Fix a vertex $v$ of $H$.  The restriction map
\[
  \cC\longrightarrow 2^{\delta_H(v)},\qquad
  A\longmapsto A\cap\delta_H(v)
\]
is injective on every connectivity code $\cC$.  Indeed, if two distinct
codewords had the same restriction, their symmetric difference would
contain no edge incident with $v$, so $v$ would be isolated.  The graph
$H$ is four-regular; hence $|\cC|\leq 2^4=16$.  Notice that this argument
does not assume linearity.

It remains to attain the bound.  Number the vertex $(i,j)$ by $4i+j$.
Number the 24 edges by
\begin{align*}
 e_{4i+j}&=\{(i,j),(i+1\bmod 3,j)\},\\
 e_{12+4i+j}&=\{(i,j),(i,j+1\bmod 4)\},
 \qquad 0\leq i<3,\quad 0\leq j<4.
\end{align*}
We identify an edge set with a 24-bit integer, with $e_r$ represented by
bit $r$ (the least significant bit is bit $0$).  In $\F_2^{24}$ set
\[
 g_0=\bits{470c7f},\qquad g_1=\bits{296bf4},\qquad
 g_2=\bits{fb304a},\qquad g_3=\bits{17e0e8},
\]
where the displayed strings are hexadecimal.  For
$a=\sum_{r=0}^3 a_r2^r\in\{0,\ldots,15\}$ define
\[
  W_a=\bigtriangleup_{r:a_r=1}g_r,
  \qquad
  \cC=\{W_a:0\leq a<16\}.
\]

Table~\ref{tab:certificate} is a directly checkable certificate.  In each
row the mask is $W_a$, and all listed edge indices have bit $1$ in that
mask.  Moreover, the edge indices are ordered so that the first edge has
two vertices and every later edge has exactly one endpoint among the
vertices already reached.  Each row therefore lists an 11-edge spanning
tree on all 12 vertices.

\begin{table}[htbp]
\centering
\caption{A spanning-tree certificate for every nonzero generated word.}
\label{tab:certificate}
\small
\setlength{\tabcolsep}{5pt}
\begin{tabular}{@{}ccl@{}}
\toprule
$a$ (hex) & $W_a$ (hex) & ordered spanning-tree edge indices \\
\midrule
$1$ & \bits{470c7f} & $0,4,16,1,5,17,2,6,18,22,3$ \\
$2$ & \bits{296bf4} & $8,4,16,19,5,7,9,21,11,13,6$ \\
$3$ & \bits{6e678b} & $0,8,19,3,7,18,14,22,17,13,21$ \\
$4$ & \bits{fb304a} & $12,1,13,16,17,19,6,3,21,22,20$ \\
$5$ & \bits{bc3c35} & $0,12,4,19,13,20,23,18,10,5,11$ \\
$6$ & \bits{d25bbe} & $8,12,4,20,23,1,7,11,22,17,14$ \\
$7$ & \bits{9557c1} & $0,8,12,16,20,23,7,18,6,10,14$ \\
$8$ & \bits{17e0e8} & $15,3,14,7,18,13,6,17,5,16,20$ \\
$9$ & \bits{50ec97} & $0,15,4,11,14,20,7,22,2,13,1$ \\
$a$ & \bits{3e8b1c} & $8,15,4,20,3,11,9,21,18,2,17$ \\
$b$ & \bits{798763} & $0,8,15,16,19,20,1,21,6,10,22$ \\
$c$ & \bits{ecd0a2} & $12,15,1,14,5,21,22,7,23,18,19$ \\
$d$ & \bits{abdcdd} & $0,12,15,4,16,19,11,14,17,10,21$ \\
$e$ & \bits{c5bb56} & $8,12,15,4,23,1,9,13,22,2,18$ \\
$f$ & \bits{82b729} & $0,8,12,15,23,9,13,3,5,10,17$ \\
\bottomrule
\end{tabular}
\end{table}

Consequently every nonzero $W_a$ is connected and nonempty.  The 16 words
are distinct: if $W_a=W_b$ for $a\neq b$, then
$W_{a\mathbin{\mathrm{XOR}}b}=0$, contrary to the table.  Finally,
\[
  W_a\mathbin{\triangle}W_b=W_{a\mathbin{\mathrm{XOR}}b},
\]
which is connected whenever $a\neq b$.  Thus $\cC$ is a 16-word
connectivity code.  Together with the upper bound, this proves
Theorem~\ref{thm:main}.

\section{Exact replay and scope}

The displayed table is sufficient for a manual finite verification, but
the construction was also replayed from serialized data.  A verifier
regenerates $C_3\mathbin{\square}C_4$ from its coordinates, reconstructs
all 16 words, checks the 15 displayed spanning trees, and independently
runs graph traversal on all $\binom{16}{2}=120$ pairwise differences.
Every test uses integer and bit operations; floating-point comparisons are
absent.  The four generators were discovered with an OR-Tools CP-SAT
model imposing local rank and integral-flow connectivity constraints, but
the solver is outside the proof's trust boundary: the theorem follows
from the generators and the spanning-tree table alone.

This separation of discovery from exact certification is also used in
recent work on finite code and clique problems, where generated objects
are reconstructed with integer predicates and checked by independent
implementations~\cite[Sections~4 and 6]{ZengEtAl2026}.  Here the finite
certificate establishes only the lower bound for the specified graph;
the unrestricted upper bound is the deductive restriction argument above.

\paragraph{Computational and AI assistance.}
OpenAI Codex assisted with the literature search, the formulation and
implementation of the discovery computation, certificate extraction,
independent exact replay, adversarial audit organization, and manuscript
preparation.  OR-Tools CP-SAT was used only to discover the four displayed
generators.  Language-model output and solver status are not mathematical
evidence for Theorem~\ref{thm:main}; its proof is the restriction argument
and the explicit spanning-tree certificate above.  No proof assistant was
used.

\bibliographystyle{plain}
\bibliography{references}

\end{document}
